% GENERATED FILE. Edit canonical lessons, then run npm run generate:books.
% canonical-source-hash: fnv1a64:696a799853061b36
% canonical-lessons: KA-C230-gere, KA-C230-hundi, KA-C230-sutkes, KA-C230-guddali, KA-C230-kudugolu

\chapter{Line, Money box, Suitcase, Hoe, Sickle}
\label{ch:ka-line-money-box-suitcase-hoe-sickle}
\hlchaptermodality{230}

\begin{chapteropening}
\textbf{By the end of this chapter:} \emph{I can say a line, a money box, a suitcase, a hoe and a sickle.}

Complete the last lesson of chapter 230: I can say a line, a money box, a suitcase, a hoe and a sickle.
\end{chapteropening}

\section[\texorpdfstring{\kn{ಗೆರೆ} (gere)}{ಗೆರೆ (gere)}]{\kn{ಗೆರೆ} (gere) — a line}

\label{lesson:KA-C230-gere}



\begin{quote}
Before the new one: say the Kannada for a knot, then the Kannada for a chain.
\end{quote}

\subsection*{You'll want to know: \kn{ಗೆರೆ}}
\textbf{\kn{ಗೆರೆ}} — \emph{gere} — \textquotedblleft{}a line\textquotedblright{}.

\begin{grammarlens}[title={What you've built}]
One more word for this chapter.
\end{grammarlens}

\subsection*{Your turn}
\begin{itemize}
\raggedright
  \item \emph{Say it:} \emph{gere}
  \item \emph{Say it:} \emph{gere}, once more
  \item \emph{Say it:} say \emph{gere}
  \item \emph{Recall:} say the Kannada for a knot, then the Kannada for a chain, then say \emph{gere} again
\end{itemize}

\begin{tcolorbox}[breakable,colback=teal!4,colframe=teal!35!black,title={Before you move on}]
What does \kn{ಗೆರೆ} mean? (\textquotedblleft{}A line\textquotedblright{}.) Say it once more.
\end{tcolorbox}
\section[\texorpdfstring{\kn{ಹುಂಡಿ} (huṇḍi)}{ಹುಂಡಿ (huṇḍi)}]{\kn{ಹುಂಡಿ} (huṇḍi) — a money box, a collection box}

\label{lesson:KA-C230-hundi}



\begin{quote}
Before the new one: say the Kannada for a chain, then the Kannada for a line.
\end{quote}

\subsection*{You'll want to know: \kn{ಹುಂಡಿ}}
\textbf{\kn{ಹುಂಡಿ}} — \emph{huṇḍi} — \textquotedblleft{}a money box, a collection box\textquotedblright{}.

\begin{grammarlens}[title={What you've built}]
One more word for this chapter.
\end{grammarlens}

\subsection*{Your turn}
\begin{itemize}
\raggedright
  \item \emph{Say it:} \emph{huṇḍi}
  \item \emph{Say it:} \emph{huṇḍi}, once more
  \item \emph{Say it:} say \emph{huṇḍi}
  \item \emph{Recall:} say the Kannada for a chain, then the Kannada for a line, then say \emph{huṇḍi} again
\end{itemize}

\begin{tcolorbox}[breakable,colback=teal!4,colframe=teal!35!black,title={Before you move on}]
What does \kn{ಹುಂಡಿ} mean? (\textquotedblleft{}A money box, a collection box\textquotedblright{}.) Say it once more.
\end{tcolorbox}
\section[\texorpdfstring{\kn{ಸೂಟ್ಕೇಸ್} (sūṭkēs)}{ಸೂಟ್ಕೇಸ್ (sūṭkēs)}]{\kn{ಸೂಟ್ಕೇಸ್} (sūṭkēs) — a suitcase}

\label{lesson:KA-C230-sutkes}



\begin{quote}
Before the new one: say the Kannada for a line, then the Kannada for a money box.
\end{quote}

\subsection*{You'll want to know: \kn{ಸೂಟ್ಕೇಸ್}}
\textbf{\kn{ಸೂಟ್ಕೇಸ್}} — \emph{sūṭkēs} — \textquotedblleft{}a suitcase\textquotedblright{}.

\begin{grammarlens}[title={What you've built}]
One more word for this chapter.
\end{grammarlens}

\subsection*{Your turn}
\begin{itemize}
\raggedright
  \item \emph{Say it:} \emph{sūṭkēs}
  \item \emph{Say it:} \emph{sūṭkēs}, once more
  \item \emph{Say it:} say \emph{sūṭkēs}
  \item \emph{Recall:} say the Kannada for a line, then the Kannada for a money box, then say \emph{sūṭkēs} again
\end{itemize}

\begin{tcolorbox}[breakable,colback=teal!4,colframe=teal!35!black,title={Before you move on}]
What does \kn{ಸೂಟ್ಕೇಸ್} mean? (\textquotedblleft{}A suitcase\textquotedblright{}.) Say it once more.
\end{tcolorbox}
\section[\texorpdfstring{\kn{ಗುದ್ದಲಿ} (guddali)}{ಗುದ್ದಲಿ (guddali)}]{\kn{ಗುದ್ದಲಿ} (guddali) — a hoe, a pickaxe}

\label{lesson:KA-C230-guddali}



\begin{quote}
Before the new one: say the Kannada for a money box, then the Kannada for a suitcase.
\end{quote}

\subsection*{You'll want to know: \kn{ಗುದ್ದಲಿ}}
\textbf{\kn{ಗುದ್ದಲಿ}} — \emph{guddali} — \textquotedblleft{}a hoe, a pickaxe\textquotedblright{}.

\begin{grammarlens}[title={What you've built}]
One more word for this chapter.
\end{grammarlens}

\subsection*{Your turn}
\begin{itemize}
\raggedright
  \item \emph{Say it:} \emph{guddali}
  \item \emph{Say it:} \emph{guddali}, once more
  \item \emph{Say it:} say \emph{guddali}
  \item \emph{Recall:} say the Kannada for a money box, then the Kannada for a suitcase, then say \emph{guddali} again
\end{itemize}

\begin{tcolorbox}[breakable,colback=teal!4,colframe=teal!35!black,title={Before you move on}]
What does \kn{ಗುದ್ದಲಿ} mean? (\textquotedblleft{}A hoe, a pickaxe\textquotedblright{}.) Say it once more.
\end{tcolorbox}
\section[\texorpdfstring{\kn{ಕುಡುಗೋಲು} (kuḍugōlu)}{ಕುಡುಗೋಲು (kuḍugōlu)}]{\kn{ಕುಡುಗೋಲು} (kuḍugōlu) — a sickle}

\label{lesson:KA-C230-kudugolu}



\begin{quote}
Before the new one: say the Kannada for a suitcase, then the Kannada for a hoe.
\end{quote}

\subsection*{You'll want to know: \kn{ಕುಡುಗೋಲು}}
\textbf{\kn{ಕುಡುಗೋಲು}} — \emph{kuḍugōlu} — \textquotedblleft{}a sickle\textquotedblright{}.

\begin{grammarlens}[title={What you've built}]
One more word for this chapter.
\end{grammarlens}

\subsection*{Your turn}
\begin{itemize}
\raggedright
  \item \emph{Say it:} \emph{kuḍugōlu}
  \item \emph{Say it:} \emph{kuḍugōlu}, once more
  \item \emph{Say it:} say \emph{kuḍugōlu}
  \item \emph{Recall:} say the Kannada for a suitcase, then the Kannada for a hoe, then say \emph{kuḍugōlu} again
\end{itemize}

\begin{tcolorbox}[breakable,colback=teal!4,colframe=teal!35!black,title={Before you move on}]
What does \kn{ಕುಡುಗೋಲು} mean? (\textquotedblleft{}A sickle\textquotedblright{}.) Say it once more.
\end{tcolorbox}
